\begin{table}[htbp]\centering\small
\caption{Irradiancia de pico frente al único umbral de daño en onda continua publicado para un sensor de silicio con su tamaño de mancha, 49\,\si{\kilo\watt\per\square\centi\meter} sobre una mancha efectiva de 9,08\,\si{\micro\meter} de radio. La hipótesis A supone que el daño lo gobierna la conducción en el silicio, de modo que el umbral reescala como $1/a$ y baja mucho para una imagen solar milimétrica; la hipótesis B supone que lo gobierna la absorción en una capa fina de baja conductividad, en cuyo caso el umbral no depende del tamaño de la mancha. A es la conservadora.}
\label{tab:threshold}
\begin{tabular}{lS[table-format=3.1]S[table-format=4.0]S[table-format=4.0]S[table-format=5.0]}
\toprule
Configuración & {$E_{\mathrm{pico}}$} & {Umbral A} & {Margen A} & {Margen B} \\
 & \multicolumn{2}{c}{(\si{\watt\per\square\centi\meter})} & & \\
\midrule
Tamron 300\,mm f/6,3, C1 del eclipse & 18.3 & 323 & 18 & 2672 \\
Tamron 16\,mm f/3,5, C1 del eclipse & 59.4 & 6059 & 102 & 825 \\
móvil f/1,9, C1 del eclipse & 201.6 & 13989 & 69 & 243 \\
Tamron 300\,mm f/6,3, mediodía & 33.7 & 323 & 10 & 1452 \\
300\,mm f/2,8, mediodía & 170.9 & 323 & 2 & 287 \\
600\,mm f/4, mediodía & 83.7 & 162 & 2 & 585 \\
\bottomrule
\end{tabular}
\end{table}
